\newcommand{\figuraVasosComunicantes}{%
\begin{figure}[htbp]
    \centering
    \begin{tikzpicture}[>=Latex,font=\small]
        \fill[cyan!20]
            (0.8,0.7)--(0.8,4.5)--(2.3,4.5)--(2.3,1.2)--
            (4.1,1.2)--(4.1,3.8)--(5.8,3.8)--(5.8,1.2)--
            (7.8,1.2)--(7.8,5.0)--(9.6,5.0)--(9.6,0.7)--cycle;
        \draw[very thick]
            (0.8,5.2)--(0.8,0.7)--(9.6,0.7)--(9.6,5.7);
        \draw[very thick] (2.3,5.2)--(2.3,1.2)--(4.1,1.2)--(4.1,4.6);
        \draw[very thick] (5.8,4.6)--(5.8,1.2)--(7.8,1.2)--(7.8,5.7);

        \draw[very thick,blue!65!black]
            (0.8,4.5)--(2.3,4.5);
        \draw[very thick,blue!65!black]
            (4.1,3.8)--(5.8,3.8);
        \draw[very thick,blue!65!black]
            (7.8,5.0)--(9.6,5.0);

        \draw[dashed,red!70!black,thick] (0.35,2.15)--(10.05,2.15);
        \foreach \x/\nombre in {1.55/A,4.95/B,8.70/C}
        {
            \fill[red!75] (\x,2.15) circle (0.10);
            \node[red!75!black,anchor=south] at (\x,2.28) {$\nombre$};
        }

        \node[font=\bfseries] at (5.2,6.05) {Vasos comunicantes};
        \node[align=center] at (5.2,0.05)
            {mismo líquido y misma altura:\quad $p_A=p_B=p_C$};
    \end{tikzpicture}
    \caption{En un líquido conectado, la presión es igual en todos los puntos
    situados a la misma altura, independientemente de la forma y del ancho de
    los recipientes.}
    \label{fig:vasos-comunicantes}
\end{figure}%
}